% Ad Atticum 13.28 (ad-atticum-13-28) --- English translation
% Translated by Claude (Anthropic) from the Latin (Perseus).
% Style guide: STYLE.md.

\ciceroLetterOpener{CICERO ATTICO salutem}

\ciceroSection{1} Since today you were to look the gardens over, what you made of them I shall of course hear tomorrow. About Faberius, whenever he comes.

\ciceroSection{2} As for the letter to Caesar --- believe me, I swear it --- I cannot. And it is not the indecency of the thing that puts me off, though indecency, more than anything else, ought to. For how disgraceful is flattery, when the very fact of being alive is disgraceful for us! But, as I began to say, this disgrace is not what holds me back. I do wish I could do it (I would then be what I ought to be); but nothing comes into my head. For consider those exhortations to Alexander composed by men of eloquence and learning --- you see what sort of ground they cover. They urge a young man on fire with the truest love of glory, eager to be given some piece of counsel that may carry him to everlasting renown, to splendour. The words are not lacking; but what can I do? Still, I had carved something or other out of an oak, something to look at least like an image. In it, because there were certain bits a little better than the things being done now or that have been done, the work was criticized; and I am not at all sorry for that. For if that letter had gone through, believe me, we would be sorry now.

\ciceroSection{3} Why, do you not see that the very pupil of Aristotle, a man of supreme talent and supreme self-restraint, became, once he was hailed as king, arrogant, cruel, ungoverned? And this messmate of our own Quirinus's procession --- do you really suppose he would be pleased by these moderate letters of ours? No, rather let him miss what is not written than disapprove of what is. In the end, let him have it as he likes. Gone is the impulse that was urging me when I gave you Archidemus's puzzle \textit{[Greek: probl\=ema Archid\=emou]}. By Hercules, I far more now wish for the outcome I then feared, or for whatever else happens. If nothing else delays you, come at a moment to suit me. Nicias, sent for with great urgency by Dolabella (for I have read the letter), has set off --- against my wishes, though on my own advice all the same.

\ciceroSection{4} This in my own hand. As I was asking Nicias about the philologists, as though it were on some other subject, we fell upon Thalna. He said nothing extreme about the man's ability --- modest, frugal. But here is what did not please me. He said he knew that Cornificia, the daughter of Quintus --- pretty long in the tooth and many times married --- had been sought by him recently in marriage; and that the women did not approve, because they were finding the estate not more than 800,000. This I thought it right for you to know.
